\section{Derivations}\label{sec:derivations}

Building on the single microstate established in \S\ref{sec:methods}, this section
presents one demonstrative line: the object generates a tower, and the tower generates
the duality web. The line runs in the order $3.1\to3.2\to3.3\to3.4\to3.5$, summarised as a
$20$-step spine at the end (\S\ref{ssec:spine}). The node that stitches \S\ref{sec:methods}
to here is the function $\Gamma$; every quantitative statement below is verified numerically
at $25$-digit precision where the identity requires it (\S\ref{ssec:status}).

% --- 3.1 ---
\subsection{The fundamental object in physics}\label{ssec:object}
The single object is the half-integer $\Gamma$ value, the carrier of the string microstate.
The seed inherited from \S\ref{sec:methods} is
\begin{equation}\label{eq:gamma-half}
  \Gamma\!\left(\tfrac12\right)=\sqrt{\pi}.
\end{equation}
The worldline phase and its modulus are
\begin{equation}\label{eq:worldline}
  \Om(\tau)=\tfrac12\,e^{\,i\tau\ln\golden},\qquad |\Om|=\tfrac12 .
\end{equation}
The wavefunction is Gaussian, with the half-factorial normalisation $\Gamma(\tfrac12)=\sqrt\pi$
and a weight fixed by the scale quantum,
\begin{equation}\label{eq:gaussian}
  \int_{\R} e^{-x^{2}}\,dx=\sqrt{\pi},\qquad
  \mu=\golden^{-\lng}=\tfrac12\quad(\golden^{\lng}=2).
\end{equation}

\begin{remark}[The wavefunction is the self-dual Gaussian]\label{rmk:wf-selfdual}
The Gaussian weight of \eqref{eq:gaussian} is the self-dual one fixed in
\S\ref{ssec:origins}, Proposition~\ref{prop:selfdual-gaussian}: normalised, and equal to its own
Fourier transform, $g(x)=e^{-\pi x^{2}}$ \eqref{eq:selfdual-gaussian}. The physical reading added
here is that this $g$ is the worldline wavefunction of the microstate; the width is not a
modelling choice but the fixed point of the Fourier transform, the same self-duality that
\S\ref{ssec:construction} used to fix $\tau=i$. \ProofWitness{selfdual\_gaussian\_unique}

\end{remark}

The one-dimensional worldline is promoted to the two-dimensional worldsheet, governed by the
Polyakov action~\cite{Polchinski98}
\begin{equation}\label{eq:polyakov}
  S_{P}=-\frac{1}{4\pi\alpha'}\int d^{2}\sigma\,\sqrt{-h}\;
        h^{ab}\,\partial_{a}X^{\mu}\,\partial_{b}X^{\nu}\,\eta_{\mu\nu}.
\end{equation}
% \HypWitness{}

\begin{proposition}[The worldsheet inherits the non-orientable fibre]\label{prop:mobius}
The worldsheet carries the non-orientable fibre of the base torus
(Proposition~\ref{prop:mobius-torus}, \S\ref{ssec:mobius-torus}): the module direction of
$\Om(\tau)=\tfrac12\,e^{\,i\tau\ln\golden}$~\eqref{eq:worldline}, carried once around the fundamental
domain, closes with the orientation-reversing half-turn $e^{i\pi}=-1$. The worldsheet is therefore
non-orientable and one-sided---the half-turn carries the module's interior edge to its exterior---so
the bulk--boundary relation the torus carries admits no privileged side, independently of which torus
carries the fibre.
\ProofWitnessBlock{fibre\_monodromy, fibre\_reflection\_order\_two, fibre\_independent\_of\_base}
\end{proposition}

The Schwinger parametrisation, used to assemble propagators and amplitudes, follows from the
$\Gamma$ integral by the change of variable $t=xA$,
\begin{equation}\label{eq:schwinger}
  A^{-n}=\frac{1}{\Gamma(n)}\int_{0}^{\infty}x^{\,n-1}\,e^{-xA}\,dx\qquad(A>0,\ n>0).
\end{equation}

\begin{remark}[The role of $\golden$ in the object]
Since $\pi=5\arccos(\golden/2)$ (\S\ref{sec:methods}), both $\sqrt\pi$ and the half-factorial
are functions of $\golden$; and $\mu=\tfrac12=\golden^{-\lng}$. The object is $\golden$-generated,
not posited.
\end{remark}

% --- 3.2 ---

\subsection{The meta-invariant: the holographic quantum}\label{ssec:meta}
The meta-invariant emerges from the structure already in place: the worldsheet phase
$\Om(\tau)=\tfrac12\,e^{\,i\tau\ln\golden}$ of \S\ref{ssec:object}~\eqref{eq:worldline},
carried on the torus $T^2_{\PCF}$ of \S\ref{ssec:construction}~\eqref{eq:torus}, whose single
modulus $|\Om|=\tfrac12$ is the norm product fixed in
\S\ref{ssec:construction}~\eqref{eq:omega-norms}. The bridge from this object to the rest of
physics is that $|\Om|=\tfrac12$ \emph{is} the holographic bit (one bit per Planck area); that this
bit is the regulator of $\Q(\sqrt5)$ measured in $\lambda_{\log}$, and not a coincidence with
$\mu_3$, is Theorem~\ref{thm:entropy-bridge}. The
spectral invariants are
\begin{equation}\label{eq:spectral-invariants}
  d=3\;(=|S_{3}|=3!=6),\qquad \mu=\tfrac12,\qquad \sigma=\tfrac32,\qquad \sigma+\mu=2 .
\end{equation}
The two $\tfrac32$'s of the corpus---the $\Gamma$ argument of \eqref{eq:half-factorial} and the
spectral exponent here---are one value for a discriminating reason: the first is $1+\mu$, the
second $3\mu$, and $1+\mu=3\mu$ has the unique solution $\mu=\tfrac12$; at arity $n$ the Basel
route \eqref{eq:sigma-basel} gives $n^{2}/6$, which meets $\tfrac32$ only at $n=3$. \ProofWitness{three\_halves\_unique, arity\_control\_three}

\paragraph{The spectral angle and the uniqueness of the invariants.}
At scale $\sigma$ the worldline subtends the spectral angle
$\alpha(\sigma)=\arctan(\epsz\golden^{\sigma})$, normalised by $\epsz M_{\PCF}=\pi$.
The invariants \eqref{eq:spectral-invariants} are not chosen but
forced: the pair $(\sigma,\mu)$ with $\sigma+\mu=2$ and $\sigma\mu=\tfrac34$ has the unique
solution $(\sigma,\mu)=(\tfrac32,\tfrac12)$,
i.e.\ $\sigma=\tfrac32$
the same arity $3$ and modulus $\tfrac12$ forced by the No-Diagonal theorem
of \S\ref{ssec:arity}. Which \emph{integer} levels the angle is read at is not settled here: it is
fixed by the level interval of \S\ref{sec:implications}, Proposition~\ref{prop:interval-uniqueness}.

\begin{proposition}[The tangent of the spectral angle is the tower]\label{prop:spectral-angle-tower}
For every $\sigma$,
\begin{equation}\label{eq:spectral-angle}
  \tan\alpha(\sigma)=\epsz\golden^{\sigma},
  \qquad\text{hence}\qquad
  \frac{\tan\alpha(\sigma+1)}{\tan\alpha(\sigma)}=\golden ,
\end{equation}
so the angle is not an auxiliary parameter: its tangent is the multiplicative spectrum that
$\golden$ generates, climbing at the one rate $\golden$ of \eqref{eq:base} --- the same spectrum
\S\ref{ssec:p-tower} will read as the tower. The surface
$\sin\alpha(\sigma_1)\cos\alpha(\sigma_2)$ has the closed form
\begin{equation}\label{eq:spectral-surface}
  \sin\alpha(\sigma_1)\cos\alpha(\sigma_2)
  =\frac{\epsz\golden^{\sigma_1}}
        {\sqrt{\bigl(1+\epsz^{2}\golden^{2\sigma_1}\bigr)\bigl(1+\epsz^{2}\golden^{2\sigma_2}\bigr)}} .
\end{equation}
What the certainty principle $\epsz\Mpcf=\pi$ \eqref{eq:certainty} contributes is the
\emph{scale} of the argument, $\epsz=\pi/\Mpcf$, and nothing else: the bound
$\alpha(\sigma)<\pi/2$ is a property of $\arctan$ and holds for every real argument, so it is
not the certainty relation that keeps the angle off the right angle. The bound is used for one
thing---it gives $\cos\alpha(\sigma)>0$, the hypothesis under which
Corollary~\ref{cor:bridge-angle} divides by $\cos\alpha$. \ProofWitness{spectral\_angle\_tan, spectral\_angle\_tower\_ratio}

\end{proposition}
\begin{proof}
$\tan\arctan x=x$ with $x=\epsz\golden^{\sigma}>0$ gives the first identity, and the quotient of two
consecutive values is $\golden^{(\sigma+1)-\sigma}=\golden$, the $\epsz$ cancelling. For
\eqref{eq:spectral-surface}, $\sin\arctan x=x/\sqrt{1+x^{2}}$ and $\cos\arctan y=1/\sqrt{1+y^{2}}$;
multiplying the two gives the display. The bound is $\arctan x<\pi/2$ for all finite $x$.
\end{proof}
With the binary entropy $H(p)=-p\log_{2}p-(1-p)\log_{2}(1-p)$, the modulus saturates the
Bekenstein bound~\cite{Bekenstein73,Hawking75},
\begin{equation}\label{eq:bridge-BH}
  \frac{1}{4\GN}=\tfrac12=|\Om|,\qquad H\!\left(\tfrac12\right)=1\ \text{bit},\qquad
  \SBH=\frac{A}{4\GN}.
\end{equation}

\begin{remark}[The role of $\golden$ in the meta-invariant]
$\mu=\tfrac12=\golden^{-\lng}$ is the scale quantum ($\golden^{\lng}=2=$ one bit). The invariants
$d=3,\ \mu=\tfrac12,\ \sigma=\tfrac32$ come from the $\golden/S_{3}$ geometry of
\S\ref{sec:methods}; they are not postulated.
\end{remark}

\paragraph{Thread (the invariants force the holographic principle).}
The holographic principle is read off the invariants, not assumed:
\begin{enumerate}
  \item the modulus $\mu=\tfrac12=|\Om|$ has binary entropy $H(\tfrac12)=1$, so each fundamental
        cell (each eigenvalue) carries exactly one bit
        (\eqref{eq:spectral-invariants}$\to$\eqref{eq:bridge-BH});
  \item the $d=3$ eigenvalues give $3$ bits $=c=3$ (Brown--Henneaux~\cite{BrownHenneaux86}: three eigenvalues $\times$
        one bit);
  \item and $1/(4\GN)=\mu=\tfrac12$ turns $\SBH=A/(4\GN)$ into the area law.
\end{enumerate}
Hence information $=$ one bit per Planck area (the holographic bound), read from the invariants.
(Numerically verified: \eqref{eq:certainty}, \eqref{eq:obs-interface}, \eqref{eq:brown-henneaux}.)

This is the holographic bound that Susskind, following 't~Hooft, drew from unproven
assumptions about the nonperturbative behaviour of string theory: one discrete degree of
freedom per Planck area~\cite{Susskind95,SusskindWitten98}. Here it is not assumed but
derived---the bit $H(\tfrac12)=1$, the central charge $c=3$, and the area law
$\SBH=A/4\GN$ are read off the single modulus $|\Om|=\tfrac12$.

% --- 3.3 ---

\subsection{The object as generator of the tower and the amplitudes}\label{ssec:p-tower}
Built on the golden tower of \S\ref{ssec:tower} (the Borger $\Lambda$-structure on
$\golden^{\sigma}$) and the mediated spectrum of \S\ref{ssec:spectrum}, the tower mediates
\emph{between} a binary register (the half-factorial $/\,\Gamma\,/\,\pi$;
$i$, $d{=}1{\to}2$) and a ternary one ($S_{3}/\golden$; $d{=}2{\to}3$), with $\mu=\tfrac12$ as
the geometric mean. The Huerta--Schreiber superpoint is purely binary ($\Z_{2}$)
\cite{huerta2018m}; PCF spans and mediates both. The relation is sharper than a contrast of
registers. The binary $\Z_{2}$ of the superpoint is not isolated in the framework: the binary
register and the turn are one object, since $2$ satisfies in $\F_5$ the equation $i$ satisfies in
$\C$ \eqref{eq:two-is-i} and is its image under the reduction at a prime over $5$
(Proposition~\ref{prop:four-cocone}). Its fourth order is therefore one of the two factors of the
conductor $20$ derived in \S\ref{ssec:spacings}, the other being the $5$ of the pentagon
(Theorem~\ref{thm:pentagon-id}), the two being coprime so that
$\operatorname{lcm}(4,5)=20$~\eqref{eq:conductor}. The $\Z_{20}$ that decides $\chi_5$
\eqref{eq:fib-criterion} therefore factors as $\Z_{4}\times\Z_{5}$, and each factor governs one
character and only one \eqref{eq:attribution}, neither determining the other
(Proposition~\ref{prop:attribution}): the binary side of the
superpoint and the pentagonal side of $\golden$ are its two components. That is the precise sense
in which PCF mediates rather than merely spans---the two registers are not juxtaposed but are the
coprime factors of a single conductor, and the mediation $\mu=\tfrac12$ of
\eqref{eq:tower-mediation} below is the geometric mean between them.
The M-theory object enters here: the HS tower
projects from the $\golden^{\sigma}$ tower. The dimensional chain is
\begin{equation}\label{eq:dim-chain}
  \R_{d=1}\xrightarrow{\;i^{2}=-1\;}\C_{d=2}
  \xrightarrow{\;\golden^{2}=\golden+1\;}\E^{3}_{d=3}.
\end{equation}
Mediation by $\mu=\tfrac12$ produces the two registers,
\begin{equation}\label{eq:tower-mediation}
  \golden^{\,\mu\lng}=\sqrt{2}\ \ (\text{binary}),\qquad
  \golden^{\,\mu\log_{\golden}3}=\sqrt{3}\ \ (\text{ternary}).
\end{equation}
\ProofWitnessBlock{M4\_sqrt2, M4\_sqrt3}
The tower and its mode count are
\begin{equation}\label{eq:tower-modes}
  z(\sigma)=\golden^{\sigma},\qquad \Nmodes(\sigma)=\big\lfloor\pi\,\golden^{\sigma}\big\rfloor .
\end{equation}
This is where the geometric encoding replaces the Fourier expansion of the string as the
device that fixes the mode content: \eqref{eq:tower-modes} counts modes by the cell capacity
$\pi$ of \eqref{eq:certainty} at each level of the tower rather than by harmonic order. The
two are not in conflict---the count is finer-grained, and it is tied through
\S\ref{ssec:adscft} to the information the horizon carries.

\begin{remark}[Fibonacci-adjacent, and where it parts company]\label{rmk:fib-adjacent}
The count \eqref{eq:tower-modes} is \emph{adjacent} to the Fibonacci sequence, not equal to it:
$\lvert\Nmodes(\sigma)-F\rvert\le1$ for the nearest Fibonacci $F$, with equality
$\Nmodes(\sigma)=F$ for $\sigma\le5$ --- the values $3,5,8,13,21,34$ --- and a genuine departure at
the tower ceiling, where $\Nmodes(6)=\lfloor\pi\golden^{6}\rfloor=56$ while the Fibonacci value is
$55$. The distinction matters: the factor $\pi$ of \eqref{eq:tower-modes} is the cell capacity of
\eqref{eq:certainty} and not a Fibonacci index, so the staircase is $\pi$-scaled and only
\emph{shadows} the recurrence of Theorem~\ref{thm:fib-min}. Reading the tower as Fibonacci would
put $55$ at the tower ceiling $\sigma=2n$, where the saturation entropy is largest. \ProofWitness{fib\_adjacent\_le\_one, Nmodes\_six\_ne\_fib} \NumericalWitness{thm:fib-min}

\end{remark}

\begin{remark}[The scale operator of the tower, and which spectrum is which]\label{rmk:mass-operator}
The tower $z=\golden^{\sigma}$ of \eqref{eq:tower-modes} defines the diagonal scale operator
\begin{equation}\label{eq:mass-operator}
  H\lvert\sigma\rangle=m_0\,\golden^{\sigma},\qquad
  \mathrm{spec}\,H=\{0\}\cup\{m_0\golden^{\sigma}\}_{\sigma\ge0},
\end{equation}
self-adjoint (real diagonal) and bounded below, with the vacuum at $0$; its gap and the reconstruction
it supports are established in \S\ref{sec:dof}. \ProofWitness{towerHam\_isHermitian, towerHam\_bounded\_below}

Two spectra must be kept apart here, and the framework supplies both from the same pair $\pi,\golden$.
The golden tower \eqref{eq:tower-modes} is the \emph{radial and scale} tower: its successive ratio
tends to $\golden$ and its physical mass reading is the Kaluza--Klein one, the interior KK mass
carrying the golden identity $\golden^{2}+\golden^{-2}-2=1$ (\S\ref{app:kk},
\eqref{eq:kk-numerator})---the gravity side. The \emph{confining} mass spectrum, on the gauge side, is
the Regge tower $\alpha'M_n^{2}=n$ of \eqref{eq:veneziano}, whose integer pole positions index the
Dirichlet series \eqref{eq:euler-product} and hence $\zeta$; its mass ratios are therefore
$M_n/M_1=\sqrt n$, not $\golden^{\sigma}$. The two are distinct spectra but not unrelated: the constant of the confining tower is the Regge slope $\alpha'$, the string tension $T=1/2\pi\alpha'$, and \S\ref{sec:dof} shows this tension welded to the golden tower level by level---so the radial and confining towers are two readings of one scale, related through the tension rather than conflated.
\end{remark}
\begin{figure}[!htbp]\centering\n\includegraphics[width=\linewidth]{fig3_N_modes.pdf}
\caption{The mode tower $\Nmodes(\sigma)=\lfloor\pi\golden^{\sigma}\rfloor$ \eqref{eq:tower-modes}:
(a) the Fibonacci-adjacent staircase $\{3,5,8,13,21,34,56,\dots\}$, equal to the Fibonacci values
for $\sigma\le5$ and departing at $\sigma=6$ where $56\neq55$
(Remark~\ref{rmk:fib-adjacent}); (b) the successive ratio
$\Nmodes(\sigma)/\Nmodes(\sigma-1)\to\golden$ alongside the Kaluza--Klein spectrum.}
\label{fig:tower-modes}
\end{figure}
The Huerta--Schreiber brane-bouquet tower~\cite{huerta2018m} maps to the PCF tower under a functor $F$,
\begin{equation}\label{eq:HS-tower}
  \R^{0|1}\to\R^{1,1|1}\to\cdots\to\R^{9,1|16}\to\R^{10,1|32}
  \ \xrightarrow{\;F\;}\ \{\golden^{\sigma}\}_{\sigma},
\end{equation}
the dynamics being presupposed on the HS side ($H=L_{0}+\bar L_{0}$) and derived on the PCF side.
Both sides are chains, so $F$ is the order-preserving level map---superpoint\,$\to$\,NS5 at
$\sigma=0,\dots,5$ on the hypercube $H_5$ (\S\ref{ssec:adscft}, App.~\ref{app:superpoint})---and
what it carries is the modulus, $|\Om|=\tfrac12$, $|\Om|^{2}=\tfrac14$, proved at the corner
(App.~\ref{app:web}).
Tree amplitudes follow from $\pi$ (binary, via $\Gamma$) and $\golden$ (ternary, via Regge): the
Veneziano amplitude \cite{Veneziano}
\begin{equation}\label{eq:veneziano}
  A_{4}=\frac{\Gamma(-\alpha's)\,\Gamma(-\alpha'u)}{\Gamma\!\big(1-\alpha'(s+u)\big)}
\end{equation}
has simple poles at $\alpha's=n$, $n\ge1$ (the Regge tower). The pole \emph{positions} are the integer
spectrum, so they index the tower's Dirichlet series; the residues are the Regge polynomials of
degree $n-1$ (carrying spin $\le n-1$), nonzero at every level, which certifies that every integer
level is populated. That Dirichlet series is $\zeta(s)$, which reorganises into the Euler product by
Proposition~\ref{prop:euler-product} of \S\ref{ssec:zeta},
The Dirichlet series whose index set the poles reproduce is \eqref{eq:euler-product}, proved in
\S\ref{ssec:zeta}; what is established here is the identification of that index set with the
Regge tower.

\begin{remark}[The tower's Dirichlet series is the Euler product of \S\ref{ssec:zeta}]\label{rmk:tower-euler}
The Euler product \eqref{eq:euler-product} of \S\ref{ssec:zeta} is here read on
the string side. The arithmetic supplies the identity; the amplitude supplies that the pole positions
$\alpha's=n$ of \eqref{eq:veneziano} are exactly its index set, with nonzero residue at every level.
The physical content is the identification of the index set, not the arithmetic identity. \ProofWitness{regge\_dirichlet\_eq\_zeta, regge\_tower\_is\_euler\_product}

\end{remark}

The pentagonal origin $\golden=2\cos(\pi/5)$ generates this structure: the Regge
tower (string side) and the Euler product (arithmetic side) are the same object seen from
two angles. The uniqueness of the ultraviolet completion selected here---that under
supersymmetry and the standard consistency conditions the Veneziano form is essentially
forced---is a recent result of Elvang and collaborators~\cite{Elvang2026}.
Loops are likewise $\golden$-organised; the one-loop partition function is finite,
\begin{equation}\label{eq:pcf-partition}
  \Zpcf(i)=\frac{e^{-3\pi/2}}{|\eta(i)|^{6}},
\end{equation}
with $\eta(i)$ the value established in \eqref{eq:eta-i}.

\begin{remark}[The one-loop constant is the lattice sum at the self-dual point]\label{rmk:eta-loop}
The $\eta(i)$ of \eqref{eq:pcf-partition} is not an independent constant: by
Remark~\ref{rmk:eta-i} of \S\ref{ssec:zeta} it is the Gaussian lattice sum evaluated at the fixed
point of $t\mapsto1/t$, $\Theta(1)=\golden^{\mu\lng}\eta(i)$ \eqref{eq:eta-i}, with
$\golden^{\mu\lng}=\sqrt2$ the binary register of Proposition~\ref{prop:mediation}. The finiteness
of \eqref{eq:pcf-partition} and the convergence of the completed zeta are therefore the same fact,
read at $\tau=i$ and at general $t$. \ProofWitness{theta\_from\_jacobi, theta\_fixed\_point\_unique}

\end{remark}

The tower carries a loop hierarchy---one loop $=$ colour, two loops $=$ generations---resting on
two proved anchors:
the colour ratio $1-\muthree^{2}=\tfrac34$, and the odd zeta values over the tower
(Theorem~\ref{thm:zeta-odd}), with Ap\'ery's $\zeta(3)$~\cite{Apery} as the two-loop constant.
What is \emph{not} derived is the assignment itself---that one loop reads as colour and two as
generations---which remains structural. \ProofWitness{colour\_ratio} The ultraviolet is regulated in two complementary
ways: the modular partition \eqref{eq:pcf-partition} (global, on the torus) and dimensional
regularisation at the vertex (local),
\begin{equation}\label{eq:gamma-pole}
  \Gamma\!\left(\tfrac{\varepsilon}{2}\right)=\frac{2}{\varepsilon}-\gamma_{E}+O(\varepsilon).
\end{equation}

\begin{remark}[The role of $\golden$ in the tower]
$\golden^{\mu\lng}=\sqrt2$ ties the binary register and $\golden^{\mu\log_{\golden}3}=\sqrt3$ the
ternary one, with $\mu=\tfrac12$ as mediator --- distinguishing PCF from the purely binary $\Z_{2}$
chain of HS.
\end{remark}

\paragraph{Thread (amplitudes and loops from $\pi$ and $\golden$ via the tower).}
\begin{enumerate}
  \item \emph{tree (Veneziano):} in \eqref{eq:veneziano} the $\Gamma$ factors carry the $\pi$/binary
        side ($\Gamma(\tfrac12)=\sqrt\pi$) and the poles at $\alpha's=n$ are the Regge tower (the
        $\golden$/ternary side);
  \item the Regge tower \emph{is} the Euler product \eqref{eq:euler-product}, linking amplitude and primes;
  \item \emph{loops:} the one-loop partition \eqref{eq:pcf-partition} is finite; the
        colour/generation hierarchy rests on two proved anchors, $1-\muthree^{2}=\tfrac34$ and the
        odd zeta values (Theorem~\ref{thm:zeta-odd}), with $\zeta(3)$ as the two-loop constant,
        while the assignment of loop order to colour and generations remains structural;
  \item the modes at each level are $\Nmodes(\sigma)=\lfloor\pi\golden^{\sigma}\rfloor$ \eqref{eq:tower-modes}.
\end{enumerate}
(Numerically verified: \eqref{lem:gamma-half}, \eqref{prop:veneziano}, \eqref{thm:basel}, \eqref{eq:pcf-partition}, \eqref{eq:tower-modes}.)

\paragraph{Thread (the two ultraviolet regulators coexist).}
\begin{enumerate}
  \item \emph{modular/$\eta$:} the partition \eqref{eq:pcf-partition} is finite because $\eta(i)$ is
        well defined --- modular invariance regulates the ultraviolet on the torus, with no divergence;
  \item \emph{dimensional/$\Gamma(\varepsilon/2)$:} \eqref{eq:gamma-pole} extracts the $2/\varepsilon$
        pole at the local vertex;
  \item they are complementary: $\eta$ handles the worldsheet (global, finite partition),
        $\Gamma(\varepsilon/2)$ the vertex (local pole); both give the same finite physics.
\end{enumerate}

% --- 3.4 ---

\subsection{All AdS/CFT from a single microstate, at each level}\label{ssec:adscft}
The whole AdS/CFT correspondence follows from one microstate (the object of \S\ref{ssec:object}),
not from an ensemble --- and it holds at each level of the tower. This per-level statement is the
dependency the web (\S\ref{ssec:web}) rests on. The single microstate is the object
$\Om$ with $|\Om|=\tfrac12$ (its realisation as a self-adjoint operator is \S\ref{sec:dof}); being
one state and not an ensemble, it resolves the
firewall ambiguity (AMPS) without averaging.
The throat and its warped metric are
\begin{equation}\label{eq:throat}
  ds^{2}=e^{2A(z)}\,\eta_{\mu\nu}\,dx^{\mu}dx^{\nu}+dz^{2},
  \qquad z(\sigma)=\golden^{\sigma}\ \text{of \eqref{eq:tower-modes}}.
\end{equation}
\paragraph{The five-dimensional throat: $\mathrm{AdS}_5$, $S^5$ and the ladder.}
The throat \eqref{eq:throat} is the $\mathrm{AdS}_5$ geometry whose curvature invariants are fixed
by the modulus: $R=-20$, $\Lambda_5=-6$, $G_{\mu\nu}=6$, saturating the
Breitenlohner--Freedman bound (Fig.~\ref{fig:ads-ladder}); computed in \S\ref{app:curvature}.
\ProofWitnessBlock{R\_scalar\_pcf, Lambda5\_pcf, G\_AB\_pcf, BF\_value}
Its companion five-sphere carries the Hopf chain $S^1\hookrightarrow S^5\to\mathbb{CP}^2$ with
$\chi(\mathbb{CP}^2)=3=n$ \ProofWitness{hopf\_latitude, hopf\_from\_clifford}
The same arity-$3$ of
\S\ref{ssec:arity}; the five-dimensional case is precisely the one requiring the Gaussian and
Eisenstein primes (\S\ref{sec:intro}). The discrete levels $\sigma=0,\dots,5$
(superpoint~$\to$~NS5, \S\ref{sssec:superpoint}) sit on the Hamming hypercube
$H_5=\{0,1\}^5$, $|H_5|=2^{5}$ \ProofWitness{hypercube\_card}
with
$\Lambda_\sigma=\golden^{\sigma}\Lambda_{\PCF}$ and the saturation entropy
$S(\sigma)=\pi\golden^{\sigma}$ \eqref{eq:tower-modes}. \ProofWitness{conductor\_attribution, four\_cocone, S\_tower}
\begin{figure}[!htbp]\centering\n\resizebox{!}{0.85\textheight}{\includegraphics{fig5_three_panel.pdf}}
\caption{$\mathrm{AdS}_5$ funnel $z=\golden^{\sigma}$ \eqref{eq:throat} $\leftrightarrow$ the
lattice tower $\Lambda_\sigma$ $\leftrightarrow$ the Hamming hypercube $H_5$ of the M-theory
ladder (superpoint~$\to$~NS5); $\golden^{\sigma}$ connects all three:
(a) the funnel with $\ell=1$, $G_N=\tfrac12$; levels bottom-to-top are
superpoint ($\sigma{=}0$), superstring ($\sigma{=}1$), het string ($\sigma{=}2$),
M2-brane ($\sigma{=}3$), M5-brane ($\sigma{=}4$), NS5-brane ($\sigma{=}5$), 10-brane ($\sigma{=}6$);
(b) the Eisenstein lattice $\Lambda_\sigma=\golden^{\sigma}\Lambda_{\rm PCF}$ in $\C$ with concentric $\golden$-rings;
(c) the Hamming cube $H_5=\{0,1\}^5$ coloured by Hamming weight $=$ H-S level.}
\label{fig:ads-ladder}
\end{figure}
The bulk--boundary map is an isometry---proved here as Lemma~\ref{lem:isometry}, and independently in the companion~\cite{HP}---used in \S\ref{sec:dof}, and it obeys the GKPW dictionary~\cite{Maldacena98,GKP98,Witten98},
\begin{equation}\label{eq:isometry}
  V:\mathcal H_{\mathrm{bulk}}\to\mathcal H_{\partial},\qquad V^{\dagger}V=\mathbf 1,\qquad
  Z_{\mathrm{grav}}\big[\phi|_{\partial}\big]
   =\big\langle e^{\int_{\partial}\phi_{0}\mathcal O}\big\rangle_{\mathrm{CFT}}.
\end{equation}
\paragraph{The isometry as an algebra--isometry (bulk--boundary) duality.}
The embedding $V$ \eqref{eq:isometry} is realised on the two faces of one modulus. The eigenvalue
triangle $\lambda_k=\tfrac12\omega^{k}$ ($\omega=e^{2\pi i/3}$) of Prop.~\ref{prop:pcf-norms}
\eqref{eq:norms-derived} is the Eisenstein lattice $\Z[\omega]$ (\emph{algebra} face: $120^\circ$,
$S_3$ symmetry, $A_2$ root system $=\mathfrak{su}(3)$)
\ProofWitnessBlock{eisenstein\_omega, Omega\_eigenvalues, a2\_minimal\_vectors\_exactly\_six};
the same modulus on the
$i/\pi$ side, where $(\tfrac12)!=\tfrac12\sqrt\pi$ fixes $\mu=\tfrac12$ (\S\ref{ssec:construction}),
is the Gauss lattice $\Z[i]$ (\emph{isometry} face: $\tau=i$, $\ell=1$). The $\golden$-mediation
$S_3\to\Z_4$ between them preserves $|\Om|=\tfrac12$ \ProofWitness{holographic\_area} and fixes $\tau=i$:
this is $V^{\dagger}V=\mathbf 1$ read on the lattice --- bulk (geometry, $\Z[i]$) and boundary
(algebra, $\Z[\omega]$) are two readings of one triangle, neither privileged, exactly as
\eqref{eq:isometry} requires.  for the two lattices and the mediation; the GKPW dictionary \ProofWitness{holographic\_area, hopf\_from\_clifford, hopf\_latitude, G\_AB\_pcf, Lambda5\_pcf, R\_scalar\_pcf, BF\_value, Omega\_eigenvalues, a2\_minimal\_vectors\_exactly\_six, eisenstein\_omega, hypercube\_card, holographic\_area, hopf\_from\_clifford, hopf\_latitude, G\_AB\_pcf, Lambda5\_pcf, R\_scalar\_pcf, BF\_value, Omega\_eigenvalues, a2\_minimal\_vectors\_exactly\_six, eisenstein\_omega, hypercube\_card}
of \eqref{eq:isometry} enters cited, ~\cite{Maldacena98,GKP98,Witten98}. \HypWitness{\cite{Maldacena98, GKP98, Witten98}}
\begin{figure}[!htbp]\centering\n\includegraphics[width=\linewidth,trim={0 5cm 0 5cm},clip]{fig4_top_down.pdf}
\caption{Isometry $\leftrightarrow$ algebra (bulk--boundary): left, the Eisenstein lattice $\Z[\omega]$
($120^\circ$, $S_3$ symmetry, $A_2=\mathfrak{su}(3)$, algebra/boundary);
right, the Gauss lattice $\Z[i]$
($90^\circ$, $\Z_4$ symmetry, $\tau=i$, $\ell=1$, isometry/bulk).
They are $\golden$-mediated ($S_3\to\Z_4$) with $|\Om|=\tfrac12$ preserved --- the
lattice realisation of $V:\mathcal H_{\mathrm{bulk}}\to\mathcal H_{\partial}$,
$V^{\dagger}V=\mathbf 1$ \eqref{eq:isometry}.}
\label{fig:isometry-algebra}
\end{figure}
Standard AdS/CFT integrates out the radial coordinate $z$ (it privileges a boundary); here
the isometry singles out neither side, so PCF keeps both bulk and boundary
(\eqref{eq:isometry}; ``forget $z$'' only on the Maldacena side).
The modulus is scale invariant, so the holographic relation holds at every level:
\begin{equation}\label{eq:tower-autosimilar}
  |\Om|_{\sigma}=\tfrac12\qquad\forall\,\sigma .
\end{equation}
The self-similarity is the physical face of the Frobenius lift \eqref{eq:frobenius-tower}: both
act on the golden monoid through the one composition law $\psi_p(\golden^{\,n})=(\golden^{\,n})^{p}$,
exponents multiplying while the modulus stands.
The gravitational data follow from the same $\mu=\tfrac12$, the $\tfrac34$ of the signature being
the $\lVert v\rVert^{2}$ of \eqref{eq:isometry-triad} read as the colour ratio
(Remark~\ref{rmk:spectral-origin}): \ProofWitness{bulkBoundary\_isometry, spectral\_uniqueness, bulkBoundary\_isometry}
\begin{equation}\label{eq:brown-henneaux}
  \SBH=\frac{1}{4\GN}=\mu=\tfrac12,\qquad c=\frac{3\ell}{2\GN}=3\quad(\ell=1,\ \GN=\tfrac12).
\end{equation}
This closes the circle to \S\ref{ssec:object}: the Newton constant is the worldline modulus
itself, $\GN=|\Om(\tau)|$ for every $\tau$ \eqref{eq:worldline}, and the central charge singles
that value out---$3\ell/(2G)=3$ iff $G=\tfrac12$.

\begin{remark}[The role of $\golden$ in AdS/CFT]
\sloppy Everything follows from the microstate $\Om$ (\S\ref{ssec:object}) and the tower
(\S\ref{ssec:p-tower}); the modulus is $\mu=\tfrac12$ at every level, and the AdS/CFT
signature $(\ell{=}1,\GN{=}\tfrac12,c{=}3,\mathrm{GKP}{=}\tfrac34)$ is exactly the microstate's
invariants.
\end{remark}

\paragraph{Thread (all AdS/CFT from one microstate).}
\begin{enumerate}
  \item the isometry $V$ with $V^{\dagger}V=\mathbf 1$ embeds bulk into boundary preserving the norm,
        so the duality holds from \emph{one} state, without averaging
        (\eqref{eq:worldline}$\to$\eqref{eq:isometry});
  \item the modulus is scale invariant, $|\Om|_{\sigma}=\tfrac12\ \forall\sigma$, so the holographic
        relation holds at \emph{each} level of the throat, not only at the boundary
        \eqref{eq:tower-autosimilar};
  \item from the same $\mu=\tfrac12$: $\SBH=1/(4\GN)=\tfrac12$ (area law) and $c=3$ (Brown--Henneaux)
        --- the gravity-side entropy and central charge come from the microstate invariant
        \eqref{eq:brown-henneaux};
  \item hence GKPW~\cite{GKP98,Witten98} and Ryu--Takayanagi~\cite{RyuTakayanagi06} are consequences of the isometric embedding of the microstate.
\end{enumerate}
Thus the entire AdS/CFT correspondence comes from one microstate, not an ensemble. (Numerically verified: $|\Om|_{\sigma}=\tfrac12\ \forall\sigma$ via \eqref{eq:certainty},
$\mathrm{GKP}=\tfrac34$, $\SBH=\mu$, $c=3$.)

Maldacena left the AdS boundary conditions unspecified---``we leave this interesting problem
for the future''~\cite{Maldacena98}. Here they are determined: the modular parameter is fixed
at $\tau=i$, the fixed point of S-duality
and the scale by the
certainty principle $\varepsilon_0 M_{\PCF}=\pi$.
\ProofWitnessBlock{certainty\_principle}

% --- 3.5 ---

\subsection{The M-theory type duality web from the microstate}\label{ssec:web}
The M-theory type duality web projects from the microstate (\S\ref{ssec:object}) through the
$\golden^{\sigma}$ tower (\S\ref{ssec:p-tower}); Witten's eight/eleven-dimensional object is not
required here. Witten \emph{posits} that object and \emph{conjectures} that the web follows from it
\cite{Witten95}; PCF starts from a \emph{demonstrated} microstate and \emph{generates} the tower, of
which the web is a projection. The web follows from the tower with HS (\S\ref{ssec:p-tower}) and from
per-level AdS/CFT (\S\ref{ssec:adscft}), both already stated. The shared signature, derived in
App.~\ref{app:web} from the web's two corners, is that the M-theory object and AdS/CFT are
projections of the microstate, carrying
\begin{equation}\label{eq:shared-signature}
  \underbrace{|\Om|=\tfrac12,\ |\Om|^{2}=\tfrac14}_{\text{Huerta--Schreiber (M-theory)}}\qquad
  \underbrace{\ell=1,\ \GN=\tfrac12,\ c=3,\ \mathrm{GKP}=\tfrac34}_{\text{Maldacena (AdS/CFT)}} ,
\end{equation}
which equal the microstate invariants $\mu=\tfrac12$. That the two corners are one and the
same object is Maldacena's conjecture that ``M/string theory on various Anti-deSitter
spacetimes is dual to various conformal field theories''~\cite{Maldacena98}---the definition of
M-theory he anticipated extending to five non-compact dimensions: here it is realised, both corners being
projections of the single microstate $|\Om|=\tfrac12$. The torus on which they meet is the
noncommutative torus that Connes, Douglas and Schwarz identify with M-theory
compactification~\cite{connes98} and that Manin attaches to the $\F_1$ program~\cite{manin04};
the non-orientable fibre it carries is independent of that algebra (Prop.~\ref{prop:mobius-torus}). The dualities that engage the tower are
\begin{align}
  \text{(S-duality)}\quad & \tau\to-\tfrac1\tau\ \Longleftrightarrow\ \golden\to-\tfrac1\golden
    =\bar\golden, \label{eq:bridge-S}\\
  \text{(T-duality)}\quad & R\to\alpha'/R, \label{eq:bridge-T}\\
  \text{(IIA strong coupling)}\quad & \sigma=10\to\sigma=11\ \ (\text{M-theory}), \label{eq:bridge-Witten}\\
  \text{(ER}=\text{EPR)}\quad & \text{entanglement}\leftrightarrow\text{geometry}, \label{eq:bridge-ER}\\
  \text{(de Sitter swampland)}\quad & \frac{|\partial_{\sigma}V|}{V}=\ln\golden, \label{eq:dS-swampland}\\
  \text{(asymptotic safety)}\quad & \text{UV fixed point}.\ProofWitness{c\_tendsto\_two} \label{eq:bridge-fixed}
\end{align}

\begin{remark}[S-duality is the involution of \S\ref{ssec:zeta}]\label{rmk:S-is-involution}
The S-duality \eqref{eq:bridge-S} is, in the Schwinger variable, the involution $t\mapsto1/t$ under
which the Gaussian lattice sum transforms with weight $\tfrac12$,
$\Theta(1/t)=\sqrt t\,\Theta(t)$ \eqref{eq:theta-lattice}, and whose unique fixed point $t=1$ is
$\tau=i$ (Proposition~\ref{prop:archimedean}). What \S\ref{ssec:zeta} proves is the transformation;
what \eqref{eq:bridge-S} adds is that the string's electric--magnetic duality and the Galois
conjugation $\golden\to-1/\golden$ are that same involution read on the coupling and on the field.
\ProofWitnessBlock{theta\_from\_jacobi, theta\_fixed\_point\_unique}
\end{remark}


\begin{remark}[The IIA$\to$11D bridge as Witten's Kaluza--Klein tower]
The bridge $\sigma=10\to\sigma=11$ in \eqref{eq:bridge-Witten} corresponds to Witten's Kaluza--Klein
mechanism: at Type~IIA strong coupling the tower of states are the Kaluza--Klein modes of
eleven-dimensional supergravity on $\R^{10}\times S^1$, of mass $\sim 1/r$~\cite{Witten95}. The
$\golden^{\sigma}$ tower plays that role here. The golden Kaluza--Klein structure
($\golden^2+\golden^{-2}-2=1$; $G_5\to G_4=\tfrac14$; the Breitenlohner--Freedman comparison;
discrete-spectrum stability) is proved in Appendix~\ref{app:kk}.
\end{remark}

\begin{remark}[The role of $\golden$ in the duality web]
The dualities are projections of the microstate's $\golden^{\sigma}$ tower, not of a posited object;
the modulus is $\mu=\tfrac12$ at every level, and S-duality is the Galois conjugation
$\golden\to-1/\golden$.
\end{remark}

\begin{remark}[the web is reconstructed from its two corners]
The web this paper uses has two corners only---Huerta--Schreiber M-theory and Maldacena AdS/CFT---and both
are shown in Appendix~\ref{app:web} to carry the microstate modulus $|\Om|=\tfrac12$; the dualities are the
symmetries of the $\golden^{\sigma}$ tower that fix it. The wider string-theory content (U-duality, $E_8$, the
five string theories) is neither used nor needed here.
\end{remark}

\paragraph{Thread (microstate $\to$ tower $\to$ web).}
\begin{enumerate}
  \item the dualities are written in $\sigma/\golden$, as projections of the tower: S as
        $\golden\to-1/\golden$ (Galois conjugate), T as $R\to\alpha'/R$, IIA$\to$11D as $\sigma=10\to11$,
        ER$=$EPR, and the de Sitter swampland $|\partial_{\sigma}V|/V=\ln\golden$
        (\eqref{eq:bridge-S}--\eqref{eq:dS-swampland});
  \item the shared signature \eqref{eq:shared-signature}: the \emph{framework colimit} proves the
        M-theory object
        (modulus $\tfrac12$) and AdS/CFT ($\ell{=}1,\GN{=}\tfrac12,c{=}3,\mathrm{GKP}{=}\tfrac34$) are
        projections with the microstate invariants ($\mu=\tfrac12$)---this is the colimit of
        \emph{frameworks}, not the Euler colimit of \S\ref{ssec:tower};
  \item hence the web is reconstructed from its two corners (App.~\ref{app:web}): both carry $\mu=\tfrac12$, and
        the dualities fix the modulus; Witten's eight-dimensional object is not required.
\end{enumerate}
Thus microstate $\to$ tower $\to$ web (HS and per-level AdS/CFT). The quantitative pieces
($\golden\to-1/\golden$, the swampland constant $\ln\golden$ (the regulator $R_K$,
Definition~\ref{def:regulator}), the T-self-dual radius, and the
shared HS/Maldacena signature) are checked numerically (\eqref{eq:obs-swampland}, \eqref{eq:obs-jacobson})
and formally in \texttt{CW6\_complete\_v2.lean} (\S\ref{ssec:status}); the framework colimit
itself---the five frameworks and the PCF core as partial spectral objects, the cocone into the
vertex $T_{\PCF}$, and its universal property---is proved in the same file
(\texttt{cocone\_property}, \texttt{colimit\_universal}, \texttt{framework\_colimit}).


The remaining duality, $\mathrm{ER}=\mathrm{EPR}$, is where the microstate supplies what
Cool Horizons left open. Maldacena and Susskind observe that the ER bridge ``may be a highly
quantum object that is yet to be independently defined'', and that they ``do not have an
independent argument for its smoothness''~\cite{MaldacenaSusskind13}. In PCF the bridge is
both defined and analytic: it is the cocycle $T(\sigma_1,\sigma_2)$ of the proposition below,
smooth in the analytic sense because its denominator $1+\varepsilon_0\golden^{\sigma}$ is strictly positive at
every level, carried, with its geometry supplied by the non-orientable fibre, on the M-theory
noncommutative torus rather than on a general quantum object.

% ====== BLOQUE MOVIDO desde §4.4 (era L1785–1825): prop:er-epr + fig + proof ======
\begin{proposition}[$\mathrm{ER}=\mathrm{EPR}$ from the modulus]\label{prop:er-epr}
The entanglement of the two hemispheres and the geometric bridge between them are one
invariant, $\mu_3^{2}=\tfrac14$.
\emph{(EPR.)} Since $|\Om(\tau)|=\tfrac12$ for every $\tau$ \eqref{eq:worldline}, the two-point
entanglement modulus is
\[
  |\Om(\tau_A)\,\overline{\Om(\tau_B)}|=|\Om(\tau_A)|\,|\Om(\tau_B)|=\mu_3^{2}=\tfrac14
  \qquad(\forall\,\tau_A,\tau_B),
\]
which is the holographic area factor $\mu_3^{2}=\tfrac14$. % FIX: era \eqref{eq:obs-cramerrao} (§4), ahora directo (autocontenido en §3)
\emph{(ER.)} The bridge amplitude between two levels of the throat,
$T(\sigma_1,\sigma_2)=\dfrac{1+\epsz\golden^{\sigma_1}}{1+\epsz\golden^{\sigma_2}}$,
satisfies $T(\sigma_1,\sigma_2)\,T(\sigma_2,\sigma_1)=1$ and
$T(\sigma_1,\sigma_2)\,T(\sigma_2,\sigma_3)=T(\sigma_1,\sigma_3)$, so the bridges form a group:
a round trip is the identity and composition is additive in $\sigma$.
\emph{($\mathrm{ER}=\mathrm{EPR}$.)} Both are fixed by the one scale-independent modulus
$\mu_3=\tfrac12$ ($|\Om|_\sigma=\tfrac12\ \forall\sigma$, \eqref{eq:tower-autosimilar}): the
entanglement between $P$ (south) and $F$ (north), mediated by $C$, \emph{is} the bridge that joins
them through the equatorial axis. Entanglement is geometry, derived from the modulus, not posited.
\ProofWitnessBlock{bridge\_cocycle, bridge\_inverse, bridge\_compose} % ADD fase 1: respaldo Lean del cociclo
\end{proposition} The bridge and its $\tfrac14$ identity are shown in Fig.~\ref{fig:er-epr}.

\begin{proof}
\emph{(EPR)} $\Om(\tau)=\tfrac12 e^{\,i\tau\ln\golden}$ has $|\Om(\tau)|=\tfrac12=\mu_3$ for all
$\tau$ \eqref{eq:worldline}; the product rule for moduli gives
$|\Om(\tau_A)\overline{\Om(\tau_B)}|=\mu_3^{2}=\tfrac14$, independent of $\tau_A,\tau_B$ because
$\mu_3$ comes from the $S_3$ arity (Proposition~\ref{prop:pcf-norms}), not from the phase.
\emph{(ER)} $T(\sigma_2,\sigma_1)$ is the reciprocal of $T(\sigma_1,\sigma_2)$, whence
$T(\sigma_1,\sigma_2)T(\sigma_2,\sigma_1)=1$; the telescoping of numerator and denominator gives
$T(\sigma_1,\sigma_2)T(\sigma_2,\sigma_3)=T(\sigma_1,\sigma_3)$.
\emph{($\mathrm{ER}=\mathrm{EPR}$)} The levels are equivalent because
$\epsz\golden^{\sigma}\!\cdot\!\Mpcf\golden^{-\sigma}=\pi$ at every $\sigma$ \eqref{eq:certainty}, so
$|\Om|_\sigma=\tfrac12$ is scale-invariant \eqref{eq:tower-autosimilar}; the entanglement invariant
$\mu_3^{2}$ and the bridge group thus share one origin.
\end{proof}

\begin{figure}[!htbp]\centering\n\includegraphics[width=\linewidth]{fig2_ER_bridge_identity.pdf}
\caption{$\mathrm{ER}=\mathrm{EPR}$ from the modulus (Prop.~\ref{prop:er-epr}): (a) the bridge
amplitude $T(\sigma_1,\sigma_2)$ as a cocycle, $T(1,4)\,T(4,7)=T(1,7)$; (b) the algebraic and
trigonometric forms agree, $T_{\mathrm{alg}}=T_{\mathrm{trig}}$, to machine precision --- which is
Corollary~\ref{cor:bridge-angle}, \eqref{eq:bridge-angle}, and not a numerical accident.}
\label{fig:er-epr}
\end{figure}

\begin{corollary}[The bridge is the spectral angle]\label{cor:bridge-angle}
With $\alpha(\sigma)$ of Proposition~\ref{prop:spectral-angle-tower},
\begin{equation}\label{eq:bridge-angle}
  T(\sigma_1,\sigma_2)
  =\frac{1+\tan\alpha(\sigma_1)}{1+\tan\alpha(\sigma_2)}
  =\frac{\sin\bigl(\alpha(\sigma_1)+\tfrac\pi4\bigr)\,\cos\alpha(\sigma_2)}
        {\sin\bigl(\alpha(\sigma_2)+\tfrac\pi4\bigr)\,\cos\alpha(\sigma_1)} .
\end{equation}
The first equality is \eqref{eq:spectral-angle}; the second is
$\sqrt2\,\sin(\alpha+\tfrac\pi4)/\cos\alpha=1+\tan\alpha$, and the $\tfrac\pi4$ enters for one
reason only: $\tan\tfrac\pi4=1$ is the $1$ that the cocycle adds to $\epsz\golden^{\sigma}$. So the
$\mathrm{ER}=\mathrm{EPR}$ bridge and the spectral angle are one object, and the trigonometric form
that Fig.~\ref{fig:er-epr} plots against the algebraic one is an identity, not a numerical
coincidence. \ProofWitness{bridge\_eq\_tan\_form, bridge\_pi\_quarter, spectral\_angle\_tan}

\end{corollary}
\begin{proof}
$\sin(\alpha+\tfrac\pi4)=(\sin\alpha+\cos\alpha)/\sqrt2$, so
$\sqrt2\sin(\alpha+\tfrac\pi4)/\cos\alpha=1+\tan\alpha$ whenever $\cos\alpha\neq0$, which holds
since $\alpha(\sigma)<\pi/2$. Forming the quotient at $\sigma_1$ and $\sigma_2$ cancels the
$\sqrt2$, and $\tan\alpha(\sigma_i)=\epsz\golden^{\sigma_i}$ returns the algebraic form of
Proposition~\ref{prop:er-epr}.
\end{proof}
% ====== fin del bloque movido ======

\begin{proposition}[Two towers, one microstate]\label{prop:two-towers}
Carrying boundary degrees of freedom into the bulk requires, in the standard treatments, a tower:
the Virasoro algebra with its central charge on the holographic side, and the chain of maximal
invariant central extensions of the superpoint, $\R^{0|1}\to\cdots\to\R^{10,1|32}$, on the M-theory
side (\S\ref{app:superpoint} below). Each is built independently and each is costly; covering what they
must cover takes a great deal of structure. Neither is contradicted here. What the colimit
\eqref{eq:shared-signature} adds is that they are two routes to one vertex, and this is visible in
four places at once.

First, \emph{strings and holography commute in $c$}. The central charge arrives twice---from the
worldsheet \eqref{eq:polyakov} and from Brown--Henneaux $c=3\ell/(2\GN)$ with the derived $\ell=1$
and $\GN=\tfrac12=|\Om|$ \eqref{eq:brown-henneaux}---and the two arrivals are the same component of
the vertex, not two values that agree. The one free parameter of the Virasoro tower is therefore the
microstate modulus: $3\ell/(2\GN)=3$ holds only at $\GN=\tfrac12$.

Second, \emph{the ladder is counted by the turn}. Its $32$ supercharges are $|H_5|=2^{5}$, and the
$2$ that counts them satisfies in $\F_5$ the equation the turn satisfies in $\C$
\eqref{eq:two-is-i}: the binary side of the ladder and the $\C$ level of the dimensional chain
\eqref{eq:dim-chain} are one generator, reduced. Field theory and group theory meet at the same
place: the two quadratic fields $\Q(i)$ and $\Q(\sqrt5)$ have coprime conductors $4$ and $5$, and
their product $20$ carries a character group whose two factors govern one character each and no more
\eqref{eq:attribution}.

Third, \emph{one recurrence covers the ladder}. The tower advances by $S(\sigma+1)=\golden S(\sigma)$
and the bulk--boundary ratio is $\pi/m_0$ at every level (\eqref{eq:intertwine} below); in any other base the
ratio would drift and no level-wise isometry could intertwine them. With no level privileged, the
many discrete steps are an index on one flow rather than a stack of separate structures---which is
why $\sigma=10\to11$ \eqref{eq:bridge-Witten} is one more level and not one more theory.

Fourth, \emph{dimension three is reached by an isometry, not by an algebra}. Frobenius's theorem
leaves only $1,2,4$ for a finite-dimensional associative real division algebra, and $3$ is not among
them~\cite{Frobenius1878}; yet the chain \eqref{eq:dim-chain} reaches $\E^{3}$. There is no
contradiction because no product in dimension three is ever asked for: what reaches it is the
norm-preserving---hence injective, hence faithful---embedding $V\colon\C\to\C^{3}$ of
Lemma~\ref{lem:isometry}. The algebraic bound and the geometric reach are different questions, and
the isometry is what carries one to the other, as the Frobenius lift \eqref{eq:frobenius-tower}
carries the same modulus on the arithmetic side.

Taken together: the complexity each tower needs on its own is the cost of not having the microstate.
for the four items; the Virasoro and brane-bouquet constructions themselves enter cited, \ProofWitness{psi\_functorial, psi\_one}
~\cite{BrownHenneaux86,Witten95}. \HypWitness{\cite{BrownHenneaux86, Witten95}}
\end{proposition}

\subsection{The demonstrative line (20-step spine, seeded by \S\ref{sec:methods})}\label{ssec:spine}
\begin{figure}[htb]\centering
\resizebox{\textwidth}{!}{%
\begin{tikzpicture}[
  font=\footnotesize,
  >={Stealth[length=4pt]},
  nd/.style={inner sep=2pt, outer sep=1pt},
  ap/.style={draw, rounded corners=2pt, fill=black!5, inner sep=3pt},
  every path/.style={->, draw=black!72, line width=0.45pt}]

% ===== 3.1 object (x=0) =====
\node[nd] (gh) at (0,5.6)  {$\Gamma(\tfrac12)=\sqrt\pi$};
\node[nd] (wl) at (0,4.4)  {$\Om=\tfrac12 e^{i\tau\ln\golden}$};
\node[nd] (gs) at (0,3.2)  {$\widehat g=g,\ g=e^{-\pi x^{2}}$};
\node[nd] (pk) at (0,2.0)  {$S_P=\!-\tfrac{1}{4\pi\alpha'}\!\int$};
\node[nd] (sw) at (0,0.6)  {$A^{-n}=\tfrac{1}{\Gamma(n)}\!\int$};

% ===== 3.2 meta (x=3.3) =====
\node[ap] (si) at (3.3,5.0) {$d{=}3,\ \mu{=}\tfrac12,\ \sigma{=}\tfrac32$};
\node[nd] (al) at (3.3,6.4) {$\tan\alpha(\sigma){=}\epsz\golden^\sigma$};
\node[nd] (bh) at (3.3,3.5) {$\tfrac{1}{4\GN}{=}\tfrac12{=}|\Om|,\ H{=}1$};

% ===== 3.3 tower + amplitudes (x=6.6) =====
\node[nd] (dc)  at (6.6,6.3)  {$\R{\to}\C{\to}E^3$};
\node[nd] (tm)  at (6.6,5.2)  {$\golden^{\mu\lng}{=}\sqrt2$};
\node[nd] (tmo) at (6.6,4.0)  {$z{=}\golden^\sigma,\ \Nmodes{=}\lfloor\pi\golden^\sigma\rfloor$};
\node[nd] (hs)  at (6.6,2.8)  {$\R^{0|1}{\to}\cdots{\to}\R^{10,1|32}$};
\node[ap] (vz)  at (6.6,1.6)  {$A_4{=}\tfrac{\Gamma\Gamma}{\Gamma}$};
\node[nd] (re)  at (6.6,0.4)  {$\sum n^{-s}{=}\zeta{=}\prod_p$};
\node[nd] (pp)  at (6.6,-0.8) {$Z{=}\tfrac{e^{-3\pi/2}}{|\eta|^{6}}$};
\node[nd] (gp)  at (6.6,-2.0) {$\Gamma(\tfrac\varepsilon2){=}\tfrac2\varepsilon{-}\gamma_E$};

% ===== 3.4 AdS/CFT (x=10.2) =====
\node[nd] (th)  at (10.2,4.0) {$ds^2{=}e^{2A}\eta{+}dz^2$};
\node[nd] (iso) at (10.2,2.7) {$V^\dagger V{=}\mathbf1$ (GKPW)};
\node[nd] (ta)  at (10.2,1.5) {$|\Om|_\sigma{=}\tfrac12\ \forall\sigma$};
\node[ap] (bhn) at (10.2,0.2) {$S_{BH}{=}\tfrac12,\ c{=}3$};
\node[nd] (adsR) at (10.2,5.3) {$R{=}{-}20,\ \Lambda_5{=}{-}6$};
\node[ap] (h5)   at (12.5,4.6) {$H_5{=}\{0,1\}^5,\ |H_5|{=}2^5$};
\node[nd] (s5)   at (12.5,3.4) {$S^1{\to}S^5{\to}\mathbb{CP}^2,\ \chi{=}3$};
\node[ap] (lat)  at (12.5,2.2) {$\mathbb{Z}[\omega]{=}A_2{=}\mathfrak{su}(3){\leftrightarrow}\mathbb{Z}[i]$};

% ===== 3.5 web (x=13.7) =====
\node[ap] (ss)  at (15.4,2.1) {$\substack{|\Om|{=}\tfrac12,\,|\Om|^2{=}\tfrac14\\ \ell{=}1,\GN{=}\tfrac12,c{=}3}$};
\node[nd] (bS)  at (18.4,4.6) {S: $\golden{\to}{-}\tfrac1\golden$};
\node[nd] (bT)  at (18.4,3.6) {T: $R{\to}\alpha'/R$};
\node[nd] (bW)  at (18.4,2.6) {IIA: $\sigma{=}10{\to}11$};
\node[nd] (bER) at (18.4,1.6) {ER$=$EPR};
\node[nd] (bsw) at (18.4,0.6) {$\tfrac{|\partial V|}{V}{=}\ln\golden$};
\node[nd] (bfx) at (18.4,-0.4){UV fixed pt};

% ===== edges =====
\draw (gh) -- (gs);
\draw (gh) to[bend left=4] (wl);
\draw (wl) -- (pk);
\draw (wl) -- (si);
\draw (gs) -- (si);
\draw (sw) to[bend right=6] (vz);
\draw (si) -- (bh);
\draw (dc) -- (tm);
\draw (tm) -- (tmo);
\draw (tmo) -- (hs);
\draw (tmo) -- (vz);
\draw (tmo) to[bend right=6] (pp);
\draw (tmo) -- (th);
\draw (vz) -- (re);
\draw (pp) -- (gp);
\draw (th) -- (iso);
\draw (al) -- (si);
\draw (th) -- (adsR);
\draw (hs) to[bend right=4] (h5);
\draw (tmo) to[bend left=4] (h5);
\draw (th) -- (s5);
\draw (iso) -- (lat);
\draw (s5) -- (lat);
\draw (iso) -- (ta);
\draw (ta) -- (bhn);
\draw (bh) to[bend right=18] (bhn);
\draw (bhn) -- (ss);
\draw (hs) to[bend left=10] (ss);
\draw (ss) -- (bS);
\draw (ss) -- (bT);
\draw (ss) -- (bW);
\draw (ss) -- (bER);
\draw (ss) -- (bsw);
\draw (ss) -- (bfx);

\end{tikzpicture}
}
\caption{Commutative spine of \S\ref{sec:derivations} (P1$\to$P20): the thirty-one nodes carrying the twenty steps, arranged by subsection (the object, \S\ref{ssec:object}; the meta-invariant, \S\ref{ssec:meta}; the tower and amplitudes, \S\ref{ssec:p-tower}; AdS/CFT per level, \S\ref{ssec:adscft}; the duality web, \S\ref{ssec:web}), with the thirty-four demonstrative arrows. Shaded nodes are the six apices, where two or more independent chains converge on the same value; the forced invariants $d{=}3,\mu{=}\tfrac12,\sigma{=}\tfrac32$, with three incoming routes, are the maximal one. Seeded by the arithmetic of $K=\Q(\sqrt5)$ in \S\ref{sec:methods}; mathematical, physical and conceptual layers coincide on one diagram.}
\label{fig:spine3}
\end{figure}
Each step uses only previous steps.

\begin{description}
  \item[\S2$\to$\S3.] P1 $\Gamma(\tfrac12)=\sqrt\pi$ \eqref{eq:gamma-half} (\S\ref{ssec:origins}; seed)
        $\to$ P2 the self-dual Gaussian \eqref{eq:selfdual-gaussian}, its Mellin the Archimedean
        factor \eqref{eq:gammaR-mellin} and its lattice sum completing $\zeta$ into
        $\xi$ \eqref{eq:places-product} (\S\ref{ssec:zeta}).
  \item[3.1 object.] P3 $\Om,|\Om|=\tfrac12$ \eqref{eq:worldline}~(\S\ref{ssec:construction}) $\to$ P4 Gaussian, $\mu=\tfrac12$
        \eqref{eq:gaussian}~(\S\ref{ssec:origins}; width fixed by Prop.~\ref{prop:selfdual-gaussian}) $\to$ P5 worldsheet/Polyakov \eqref{eq:polyakov}.
  \item[3.2 meta-invariant.] P6 invariants $d{=}3,\mu{=}\tfrac12,\sigma{=}\tfrac32$
        \eqref{eq:spectral-invariants}~(\S\ref{ssec:arity},\,\S\ref{ssec:commute}) $\to$ P7 $1/(4\GN)=\tfrac12=|\Om|$, $H(\tfrac12)=1$
        \eqref{eq:bridge-BH}~(\S\ref{ssec:construction}).
  \item[3.3 tower $+$ HS.] P8 chain $i\to\golden$ \eqref{eq:dim-chain}~(\S\ref{ssec:spectrum}) $\to$ P9 $\sqrt2/\sqrt3$ by $\mu$
        \eqref{eq:tower-mediation}~(\S\ref{ssec:origins}) $\to$ P10 tower $\golden^{\sigma}$, $\Nmodes$ \eqref{eq:tower-modes}~(\S\ref{ssec:tower})
        $\to$ P11 HS$\leftrightarrow$tower \eqref{eq:HS-tower} $\to$ P12 Veneziano \eqref{eq:veneziano}
        $\to$ P13 loops $\Zpcf(i)$ \eqref{eq:pcf-partition} $\to$ P14 both ultraviolet regulators \eqref{eq:gamma-pole}.
  \item[3.4 AdS/CFT per level.] P15 throat $z(\sigma)=\golden^{\sigma}$ \eqref{eq:throat}~(\S\ref{ssec:tower}) $\to$ P16
        isometry $V^{\dagger}V=\mathbf 1$, GKPW \eqref{eq:isometry} $\to$ P17 $|\Om|_{\sigma}=\tfrac12\ \forall\sigma$
        \eqref{eq:tower-autosimilar}~(\S\ref{ssec:construction}) $\to$ P18 $\SBH=\mu$, $c=3$ \eqref{eq:brown-henneaux}~(\S\ref{ssec:commute}).
  \item[3.5 web (final).] P19 shared signature \eqref{eq:shared-signature} $\to$ P20 web
        \eqref{eq:bridge-S}--\eqref{eq:bridge-fixed}, whose S-duality is the involution that
        P2 already carries (Remark~\ref{rmk:S-is-involution}).
\end{description}
\noindent The commutative diagram supports this line with precise structure. The thirty-one
nodes form a directed acyclic graph with four roots---the seed $\Gamma(\tfrac12)=\sqrt\pi$
\eqref{eq:gamma-half}, the Schwinger parametrisation \eqref{eq:schwinger}, the spectral angle
\eqref{eq:spectral-angle} and the dimensional chain \eqref{eq:dim-chain}---from which all else
derives, with $34$ demonstrative edges. A topological order exists---the steps can be stated so
that none depends on a later one---which guarantees the line reads without forward reference. Six
nodes are apices, where two or more independent chains converge; the forced invariants
$d{=}3,\mu{=}\tfrac12,\sigma{=}\tfrac32$, with three routes, are the maximal one.
Those apices are the sense in which the construction is \emph{overdetermined}: the invariants
are not posited and then shown consistent, but forced by more independent conditions than
there are degrees of freedom to satisfy them. This is what anchors the proposal in what is
already observed---it brings holography to de~Sitter by overdetermination rather than by
postulate.

\noindent The observer steps that close the introduction's problem (the objectivity threshold
$f_{\mathrm{crit}}=\tfrac12=\mu$ and the Cram\'er--Rao saturation $F_{\max}^{-1}=\tfrac14=\mu^{2}$) open
\S4 Implications. Here P19 is what makes P20 (the web) defensible without re-deriving M-theory or
positing Witten's object: \S\ref{sec:derivations} closes with its demonstrative line complete (Fig.~\ref{fig:spine3}).
% ===================================================================

\FloatBarrier
